\documentclass[10pt,a4paper]{amsart}
\usepackage{amsmath,amssymb,amsthm,graphicx}
% Retained preamble for the Stacks Project chapter snapshot in this directory.
% The chapter's own preamble.tex is a copy of the shared Stacks Project preamble; the block
% below is sliced MECHANICALLY, line for line, from it. Nothing was added, removed or
% redefined; the wrapper class, the xr-hyper/hyperref wiring and the externaldocument list
% are omitted because the assembler supplies the document class.
%
% lines taken from preamble.tex: 3, 5, 9, 11, 12, 156, 157, 158, ... (51 lines total)
%
% package coverage restored from the chapter's own preamble (2026-10-13): xy, xy-2cell

\usepackage[all]{xy}
\xyoption{2cell}
\UseAllTwocells

\newenvironment{reference}{\comment}{\endcomment}
\newenvironment{slogan}{\comment}{\endcomment}
\newenvironment{history}{\comment}{\endcomment}
\theoremstyle{plain}
\newtheorem{theorem}[subsection]{Theorem}
\newtheorem{proposition}[subsection]{Proposition}
\newtheorem{lemma}[subsection]{Lemma}
\theoremstyle{definition}
\newtheorem{definition}[subsection]{Definition}
\newtheorem{example}[subsection]{Example}
\newtheorem{exercise}[subsection]{Exercise}
\newtheorem{situation}[subsection]{Situation}
\theoremstyle{remark}
\newtheorem{remark}[subsection]{Remark}
\newtheorem{remarks}[subsection]{Remarks}
\numberwithin{equation}{subsection}
\def\lim{\mathop{\mathrm{lim}}\nolimits}
\def\colim{\mathop{\mathrm{colim}}\nolimits}
\def\Spec{\mathop{\mathrm{Spec}}}
\def\Hom{\mathop{\mathrm{Hom}}\nolimits}
\def\Ext{\mathop{\mathrm{Ext}}\nolimits}
\def\SheafHom{\mathop{\mathcal{H}\!\mathit{om}}\nolimits}
\def\SheafExt{\mathop{\mathcal{E}\!\mathit{xt}}\nolimits}
\def\Sch{\mathit{Sch}}
\def\Mor{\mathop{\mathrm{Mor}}\nolimits}
\def\Ob{\mathop{\mathrm{Ob}}\nolimits}
\def\Sh{\mathop{\mathit{Sh}}\nolimits}
\def\NL{\mathop{N\!L}\nolimits}
\def\CH{\mathop{\mathrm{CH}}\nolimits}
\def\proetale{{pro\text{-}\acute{e}tale}}
\def\etale{{\acute{e}tale}}
\def\QCoh{\mathit{QCoh}}
\def\Ker{\mathop{\mathrm{Ker}}}
\def\Im{\mathop{\mathrm{Im}}}
\def\Coker{\mathop{\mathrm{Coker}}}
\def\Coim{\mathop{\mathrm{Coim}}}
\DeclareMathSymbol{\boxtimes}{\mathbin}{AMSa}{"02}
\def\QCohstack{\mathcal{QC}\!\mathit{oh}}
\def\Cohstack{\mathcal{C}\!\mathit{oh}}
\def\Spacesstack{\mathcal{S}\!\mathit{paces}}
\def\Quotfunctor{\mathrm{Quot}}
\def\Hilbfunctor{\mathrm{Hilb}}
\def\Curvesstack{\mathcal{C}\!\mathit{urves}}
\def\Polarizedstack{\mathcal{P}\!\mathit{olarized}}
\def\Complexesstack{\mathcal{C}\!\mathit{omplexes}}
\def\Pic{\mathop{\mathrm{Pic}}\nolimits}
\def\Picardstack{\mathcal{P}\!\mathit{ic}}
\def\Picardfunctor{\mathrm{Pic}}
\def\Deformationcategory{\mathcal{D}\!\mathit{ef}}
\usepackage{cleveref}
\crefname{originallemma}{original lemma}{original lemmas}
\Crefname{originallemma}{Original lemma}{Original lemmas}
\makeatletter
\newlabel{example-representations}{{0D3J}{source}{}{sourcepointer-1}{}}
\newlabel{example-representations@cref}{{[originallemma][0][]0D3J}{[1][0][]source}{}{}{}}
\makeatother
\crefname{originallemma}{original lemma}{original lemmas}
\Crefname{originallemma}{Original lemma}{Original lemmas}
\makeatletter
\newlabel{formal-defos-lemma-RS-associated-functor}{{06J8}{source}{}{sourcepointer-2}{}}
\newlabel{formal-defos-lemma-RS-associated-functor@cref}{{[originallemma][0][]06J8}{[1][0][]source}{}{}{}}
\makeatother
\crefname{originallemma}{original lemma}{original lemmas}
\Crefname{originallemma}{Original lemma}{Original lemmas}
\makeatletter
\newlabel{formal-defos-lemma-RS-implies-S1-S2}{{06J7}{source}{}{sourcepointer-3}{}}
\newlabel{formal-defos-lemma-RS-implies-S1-S2@cref}{{[originallemma][0][]06J7}{[1][0][]source}{}{}{}}
\makeatother
\crefname{originallemma}{original lemma}{original lemmas}
\Crefname{originallemma}{Original lemma}{Original lemmas}
\makeatletter
\newlabel{formal-defos-remark-compose-minimal-into-iso-classes}{{06IZ}{source}{}{sourcepointer-4}{}}
\newlabel{formal-defos-remark-compose-minimal-into-iso-classes@cref}{{[originallemma][0][]06IZ}{[1][0][]source}{}{}{}}
\makeatother
\crefname{originallemma}{original lemma}{original lemmas}
\Crefname{originallemma}{Original lemma}{Original lemmas}
\makeatletter
\newlabel{formal-defos-theorem-Schlessinger-prorepresentability}{{06JM}{source}{}{sourcepointer-5}{}}
\newlabel{formal-defos-theorem-Schlessinger-prorepresentability@cref}{{[originallemma][0][]06JM}{[1][0][]source}{}{}{}}
\makeatother
\crefname{originallemma}{original lemma}{original lemmas}
\Crefname{originallemma}{Original lemma}{Original lemmas}
\makeatletter
\newlabel{lemma-representations-RS}{{0DVS}{source}{}{sourcepointer-6}{}}
\newlabel{lemma-representations-RS@cref}{{[originallemma][0][]0DVS}{[1][0][]source}{}{}{}}
\makeatother
\crefname{originallemma}{original lemma}{original lemmas}
\Crefname{originallemma}{Original lemma}{Original lemmas}
\makeatletter
\newlabel{lemma-representations-TI}{{0DVT}{source}{}{sourcepointer-7}{}}
\newlabel{lemma-representations-TI@cref}{{[originallemma][0][]0DVT}{[1][0][]source}{}{}{}}
\makeatother
\title{Stacks, Examples of Deformations: lemma-representations-hull}
\author{The Stacks Project Authors}
\date{stacks-project:examples-defos.tex:local-snapshot-20261012; GNU FDL 1.2 or later}
\begin{document}
\maketitle
\noindent\textit{Editorial note.} Statement span is lines 635-648 of sources/stacks-examples-defos/pinned-original-source.tex; proof span is lines 650-687, the complete author proof (36 body lines).\par\medskip
\noindent\textit{Editorial note.} Navigation stubs: 7 cited-but-not-included labels are supplied as non-load-bearing pointers with their pinned tags where the snapshot tag index has them.\par\medskip
\noindent\textit{Editorial note.} extra\_packages is empty on purpose: the retained preamble already loads xy with options and a second declaration clashes in the assembler wrapper.\par\medskip
\noindent\textit{Editorial note.} Text is retained verbatim, including the author's own typos and placeholders.\par\medskip
\noindent\textit{Editorial note.} Queue index 13 of queue\_examples-defos.json (pb=1589, stmt 635-648, proof 650-687).\par\medskip
\noindent\textit{Allowed standard prerequisites (editorial).}\begin{itemize}
\item Explicit prerequisite boundary. The cited inputs are prerequisites; the argument in the retained span is not replaced by any cited result.
\end{itemize}\medskip
\ifdefined\hypertarget\hypertarget{sourcepointer-1}{}\fi
\noindent\textit{Source reference.} Original lemma 0D3J: sources/stacks-examples-defos/pinned-original-source.tex, source lines 483--494. Cited by the retained statement or proof but not part of the collected spans; supplied only as a navigation pointer. The argument does not depend on its content.\par\medskip
\ifdefined\hypertarget\hypertarget{sourcepointer-2}{}\fi
\noindent\textit{Source reference.} Original lemma 06J8: sources/stacks-examples-defos/pinned-original-source.tex, source lines 687--687. Cited by the retained statement or proof but not part of the collected spans; supplied only as a navigation pointer. The argument does not depend on its content.\par\medskip
\ifdefined\hypertarget\hypertarget{sourcepointer-3}{}\fi
\noindent\textit{Source reference.} Original lemma 06J7: sources/stacks-examples-defos/pinned-original-source.tex, source lines 687--687. Cited by the retained statement or proof but not part of the collected spans; supplied only as a navigation pointer. The argument does not depend on its content.\par\medskip
\ifdefined\hypertarget\hypertarget{sourcepointer-4}{}\fi
\noindent\textit{Source reference.} Original lemma 06IZ: sources/stacks-examples-defos/pinned-original-source.tex, source lines 687--687. Cited by the retained statement or proof but not part of the collected spans; supplied only as a navigation pointer. The argument does not depend on its content.\par\medskip
\ifdefined\hypertarget\hypertarget{sourcepointer-5}{}\fi
\noindent\textit{Source reference.} Original lemma 06JM: sources/stacks-examples-defos/pinned-original-source.tex, source lines 687--687. Cited by the retained statement or proof but not part of the collected spans; supplied only as a navigation pointer. The argument does not depend on its content.\par\medskip
\ifdefined\hypertarget\hypertarget{sourcepointer-6}{}\fi
\noindent\textit{Source reference.} Original lemma 0DVS: sources/stacks-examples-defos/pinned-original-source.tex, source lines 526--532. Cited by the retained statement or proof but not part of the collected spans; supplied only as a navigation pointer. The argument does not depend on its content.\par\medskip
\ifdefined\hypertarget\hypertarget{sourcepointer-7}{}\fi
\noindent\textit{Source reference.} Original lemma 0DVT: sources/stacks-examples-defos/pinned-original-source.tex, source lines 564--578. Cited by the retained statement or proof but not part of the collected spans; supplied only as a navigation pointer. The argument does not depend on its content.\par\medskip
\begin{lemma}
\label{lemma-representations-hull}
In Example \ref{example-representations} assume $\Gamma$ finitely generated.
Let $\rho_0 : \Gamma \to \text{GL}_k(V)$ be a finite dimensional representation.
Assume $\Lambda$ is a complete local ring with residue field $k$
(the classical case). Then the functor
$$
F : \mathcal{C}_\Lambda \longrightarrow \textit{Sets},\quad
A \longmapsto \Ob(\Deformationcategory_{V, \rho_0}(A))/\cong
$$
of isomorphism classes of objects has a hull. If
$H^0(\Gamma, \text{End}_k(V)) = k$, then $F$ is
prorepresentable.
\end{lemma}

\begin{proof}
The existence of a hull follows from Lemmas \ref{lemma-representations-RS} and
\ref{lemma-representations-TI} and
Formal Deformation Theory, Lemma \ref{formal-defos-lemma-RS-implies-S1-S2}
and Remark \ref{formal-defos-remark-compose-minimal-into-iso-classes}.

\medskip\noindent
Assume $H^0(\Gamma, \text{End}_k(V)) = k$. To see that $F$
is prorepresentable it suffices to show that $F$ is a
deformation functor, see Formal Deformation Theory, Theorem
\ref{formal-defos-theorem-Schlessinger-prorepresentability}.
In other words, we have to show $F$ satisfies (RS).
For this we can use the criterion of Formal Deformation Theory, Lemma
\ref{formal-defos-lemma-RS-associated-functor}.
The required surjectivity of automorphism groups will follow if we
show that
$$
A \cdot \text{id}_M =
\text{End}_{A[\Gamma]}(M)
$$
for any object $(A, M, \rho)$ of $\mathcal{F}$ such that
$M \otimes_A k$ is isomorphic to $V$ as a representation of $\Gamma$.
Since the left hand side is contained in the right hand side,
it suffices to show
$\text{length}_A \text{End}_{A[\Gamma]}(M) \leq \text{length}_A A$.
Choose pairwise distinct ideals
$(0) = I_n \subset \ldots \subset I_1 \subset A$
with $n = \text{length}(A)$. By correspondingly filtering
$M$, we see that it suffices to prove $\Hom_{A[\Gamma]}(M, I_tM/I_{t + 1}M)$
has length $1$. Since $I_tM/I_{t + 1}M \cong M \otimes_A k$
and since any $A[\Gamma]$-module map $M \to M \otimes_A k$ factors
uniquely through the quotient map $M \to M \otimes_A k$
to give an element of
$$
\text{End}_{A[\Gamma]}(M \otimes_A k) = \text{End}_{k[\Gamma]}(V) = k
$$
we conclude.
\end{proof}

\end{document}
